% Options for packages loaded elsewhere
\PassOptionsToPackage{unicode}{hyperref}
\PassOptionsToPackage{hyphens}{url}
\documentclass[
  ukrainian,
]{article}
\usepackage{xcolor}
\usepackage[margin=2.5cm]{geometry}
\usepackage{amsmath,amssymb}
\setcounter{secnumdepth}{-\maxdimen} % remove section numbering
\usepackage{iftex}
\ifPDFTeX
  \usepackage[T1]{fontenc}
  \usepackage[utf8]{inputenc}
  \usepackage{textcomp} % provide euro and other symbols
\else % if luatex or xetex
  \usepackage{unicode-math} % this also loads fontspec
  \defaultfontfeatures{Scale=MatchLowercase}
  \defaultfontfeatures[\rmfamily]{Ligatures=TeX,Scale=1}
\fi
\usepackage{lmodern}
\ifPDFTeX\else
\setmainfont{DejaVu Serif}
\setsansfont{DejaVu Sans}
\setmonofont{DejaVu Sans Mono}
\setmathfont{STIX Two Math}
\fi
\ifPDFTeX\else
  % xetex/luatex font selection
\fi
% Use upquote if available, for straight quotes in verbatim environments
\IfFileExists{upquote.sty}{\usepackage{upquote}}{}
\IfFileExists{microtype.sty}{% use microtype if available
  \usepackage[]{microtype}
  \UseMicrotypeSet[protrusion]{basicmath} % disable protrusion for tt fonts
}{}
\makeatletter
\@ifundefined{KOMAClassName}{% if non-KOMA class
  \IfFileExists{parskip.sty}{%
    \usepackage{parskip}
  }{% else
    \setlength{\parindent}{0pt}
    \setlength{\parskip}{6pt plus 2pt minus 1pt}}
}{% if KOMA class
  \KOMAoptions{parskip=half}}
\makeatother
\usepackage{graphicx}
\graphicspath{{/Users/kalinicenkovitalijmikolajovic/.codex/worktrees/5c15/Script_for_translate_book/Output_Final/}}
\setkeys{Gin}{width=0.82\linewidth,height=0.82\textheight,keepaspectratio}
\makeatletter
\newsavebox\pandoc@box
\newcommand*\pandocbounded[1]{% scales image to fit in text height/width
  \sbox\pandoc@box{#1}%
  \Gscale@div\@tempa{\textheight}{\dimexpr\ht\pandoc@box+\dp\pandoc@box\relax}%
  \Gscale@div\@tempb{\linewidth}{\wd\pandoc@box}%
  \ifdim\@tempb\p@<\@tempa\p@\let\@tempa\@tempb\fi% select the smaller of both
  \ifdim\@tempa\p@<\p@\scalebox{\@tempa}{\usebox\pandoc@box}%
  \else\usebox{\pandoc@box}%
  \fi%
}
% Set default figure placement to htbp
\def\fps@figure{htbp}
\makeatother
\ifLuaTeX
\usepackage[bidi=basic,shorthands=off,provide=*]{babel}
\else
\usepackage[bidi=default,shorthands=off,provide=*]{babel}
\fi
\ifLuaTeX
  \usepackage{selnolig} % disable illegal ligatures
\fi
\setlength{\emergencystretch}{3em} % prevent overfull lines
\providecommand{\tightlist}{%
  \setlength{\itemsep}{0pt}\setlength{\parskip}{0pt}}
\usepackage{bookmark}
\IfFileExists{xurl.sty}{\usepackage{xurl}}{} % add URL line breaks if available
\urlstyle{same}
\hypersetup{
  pdflang={uk},
  hidelinks,
  pdfcreator={LaTeX via pandoc}}

\author{}
\date{}

\begin{document}

\section{ФОРМУЛА ДАЛАМБЕРА МЕТА ЛЕКЦІЇ: Построить решение задачи Коши
для волнового
уравнения}\label{ux444ux43eux440ux43cux443ux43bux430-ux434ux430ux43bux430ux43cux431ux435ux440ux430-ux43cux435ux442ux430-ux43bux435ux43aux446ux456ux457-ux43fux43eux441ux442ux440ux43eux438ux442ux44c-ux440ux435ux448ux435ux43dux438ux435-ux437ux430ux434ux430ux447ux438-ux43aux43eux448ux438-ux434ux43bux44f-ux432ux43eux43bux43dux43eux432ux43eux433ux43e-ux443ux440ux430ux432ux43dux435ux43dux438ux44f}

и \emph{\{yy\}=\(u_{l+1,\ j}\)-2u}\{l,~j\}+\(u_{l-1,\ j}\), (центральные
разностные производные),

і підставляють їх у рівняння похідних похідних. Відповідна сітка
показана на рис. 43.2. Розв'язавши отримане рівняння за \(u_{i,l}\) ,
отримуємо

Посмотрим внимательно на это уравнение. Коэффициенты при
\(u_{i+1,\ j},\ u_i,\ j+i,\ u_l,\ j-1\) н \(u_{i-1},\ j\) положительны,
и их сумма равна единице. Другими словами, решение \(u_i,\ j\) является
взвешенным средним решений в четырех соседних точках. Следовательно,
можно Р И С

\pandocbounded{\includegraphics[keepaspectratio]{images/_page_324_Picture_9.jpeg}}

. 43.2. Узлы сетки для случайного маршруга.

модифицировать «Блуждающего пьяницу» так, чтобы вероятности перехода в
соседние точки были не 0,25, а равнялись бы коэффициенту в
соответствующем члене. Другими словами, если пьяница находится в точке
(i, j), то он переходит в точку

Ця задача описує рух нескінченної струни за заданих початкових умов.
Його розв'язав у 1750 році французький математик д'Аламбер. Формула
Д'Аламбера

с вероятностью \(\frac{1}{2(1+\sin x_j)}\) , \((i, j-1)\) с вероятностью
\(\frac{1}{2(1+\sin x_j)}\) , \((i+1, j)\) с вероятностью
\(\frac{\sin x_j}{2(1+\sin x_j)}\) , \((i-1, j)\) с вероятностью
\(\frac{\sin x_j}{2(1+\sin x_j)}\) .

В остальном игра не меняется. Модификация для других задач может
оказаться более хитроумной, но основная идея остается той же. Читатель
может заняться разработкой своих игр для решения других задач.
Параболический случай рассмотрен в разделе «Задачи». \# ЗАУВАЖЕННЯ 1.
Отметим, что если относительные частоты попадания \(P_A(p_l)\) известны,
то можно легко найти величину \(u\left(A\right)\) при других граничных
условиях \(g_l\) . Для этого достаточно подставить \(g_l\) в формулу

полегшує пошук розв'язку, якщо початкові умови відомі. Крім того, це
дозволяє дати цікаву фізичну інтерпретацію розв'язку мовою хвиль, що
рухаються.

\begin{verbatim}
Як читач, ймовірно, пам'ятає, у параболічному випадку ми спочатку розв'язали задачі дифузії на скінченному відрізку (методом розділення змінних), а лише потім переходили до необмежених інтервалів (-  \infty < x < \infty  ), де застосовували перетворення Фур'є. У гіперболічному випадку ми зробимо навпаки. Почнемо з розв'язання одномірного хвильового рівняння на всій дійсній осі, тобто з задачі Коші
\end{verbatim}

(17.1)

а не разыгрывать вновь случайные прогулки.

\begin{enumerate}
\def\labelenumi{\arabic{enumi}.}
\setcounter{enumi}{1}
\tightlist
\item
  Во многих случаях исследователя интересует решение уравнения с
  частными производными в одной точке. Если граница области сложна и
  уравнение зависит от трех или четырех независимых переменных, метод
  Монте-Карло может оказаться удобным. Методы Монте-Карло возникли при
  решении очень трудных задач диффузии нейтронов. Эти задачи невозможно
  решать аналитически. \# ЗАДАЧІ
\item
  Построить блок-схему игры «Блуждающий пьяница» для решения задачи
\end{enumerate}

(УЧП)

, 0 \textless{} x \textless{} 1 , 0 \textless{} y \textless{} 1 , (ГУ)
\(u(x, y)\) = \(g(x, y)\) на границе,

\begin{itemize}
\item
  у вузлах внутрішньої сітки, якщо кількість вузлів сітки довільна. 2.
  Напишіть програму для розв'язання задачі відповідно до блок-схеми з
  завдання 1.
\item
  \begin{enumerate}
  \def\labelenumi{\arabic{enumi}.}
  \setcounter{enumi}{2}
  \tightlist
  \item
    Як модифікувати гру для вирішення проблеми
  \end{enumerate}
\end{itemize}

(УЧП)

содержащей только начальные условия. Эту задачу можно было бы решить
либо с помощью преобразования Фурье (по переменной t), либо с помощью
преобразования Лапласа (по переменной t). Однако мы воспользуемся другим
методом (методом канонических координат), который, как мы надеемся,
заинтересует читателя. Этот метод базируется на тех же идеях, что и
метод перехода к движущейся системе отсчета, с которым мы познакомились
в лекции 15. Итак, приступим к решению задачи (17.1).

\section{Решение одномерного волнового уравнения. Формула
Даламбера}\label{ux440ux435ux448ux435ux43dux438ux435-ux43eux434ux43dux43eux43cux435ux440ux43dux43eux433ux43e-ux432ux43eux43bux43dux43eux432ux43eux433ux43e-ux443ux440ux430ux432ux43dux435ux43dux438ux44f.-ux444ux43eux440ux43cux443ux43bux430-ux434ux430ux43bux430ux43cux431ux435ux440ux430}

Решение задачи (17.1) разобьем на несколько шагов:

ШАГ 1 (Замена координат (x, t) новыми каноническими координатами
\((\xi, \eta)\) ).

Для решения задачи (17.1) воспользуемся тем, что если заменить две
независимые переменные x и t новыми пространственновременными
координатами

, 0 \textless{} x \textless{} 1 , 0 \textless{} y \textless{} 1 , (ГУ)
\(u(x, y)\) = \(g(x, y)\) на границе?

Как выглядит в этой игре случайная прогулка?

\begin{enumerate}
\def\labelenumi{\arabic{enumi}.}
\setcounter{enumi}{3}
\tightlist
\item
  Можете ли вы так модифицировать игру, чтобы решить задачу
\end{enumerate}

\[
\begin{array}{l} u_{xx} + u_{yy} = 0, \quad 0 < x < 1, \quad 0 < y < 1, \\ u(x, 1) = 0, \quad 0 < x < 1, \quad 0 < y < 1, \\ u(x, 0) = 0, \quad 0 < x < 1, \quad 0 < y < 1, \\ u(0, y) = 1, \quad \frac{\partial u}{\partial x}(1, y) = 0 \end{array}
\]

то уравнение

\begin{verbatim}
\begin{aligned}
(\text{УЧП})\quad & u_t = \alpha^2 u_{xx} - \beta u_x - \gamma u, && 0 < x < 1,\ 0 < t < \infty, \\
(\text{ГУ})\quad & u(0,t) = f(t),\quad u(1,t) = g(t), && 0 < t < \infty, \\
(\text{НУ})\quad & u(x,0) = \varphi(x), && 0 \leqslant x \leqslant 1.
\end{aligned}
\end{verbatim}

преобразуется в уравнение

ПРИМІТКА. Замінимо це рівняння на скінченну різницю Кренка--Ніколсона,
виразимо \(u_{i+1,j}\) через п'ять значень \(u_{i+1,j-1}\) ,
\(u_{i+1,j+1}\) , \(u_{i,j+1}\) , \(u_{i,j+1}\) , \(u_{i,j+1}\) у
суміжних точках.

\section{КАЛЬКУЛЮС ВАРІАЦІЙ (РІВНЯННЯ
ЕЙЛЕРА--ЛАГРАНЖА)}\label{ux43aux430ux43bux44cux43aux443ux43bux44eux441-ux432ux430ux440ux456ux430ux446ux456ux439-ux440ux456ux432ux43dux44fux43dux43dux44f-ux435ux439ux43bux435ux440ux430ux43bux430ux433ux440ux430ux43dux436ux430}

МЕТА ЛЕКЦІЇ: Ввести поняття функціоналу (функцію з функції) і пояснити,
як функціонали природно виникають у фізиці. Типовий функціонал є
інтегралом

.

Це легко перевірити, якщо ви використовуєте формули перетворення
координат:

(17.2)

що є функцією функції y (припускається, що інтегративна функція
\(F\left(x,\,y,\,y'\right)\) задана). Приклад функціональності:

Подставляя эти выражения для производных в волновое уравнение, получим

Ми покажемо, як знайти функцію \(\overline{y}(x)\) , яка мінімізує
функціональність J{[}y{]}. Виявляється, що мінімізуюча функція має
\(\overline{y}\) задовольняти так зване рівняння Ейлера-Лагранжа. Це
рівняння виконує ту ж роль, що й необхідна умова мінімуму

.. Первый шаг сделан.

КРОК 2 (Розв'язання трансформованого рівняння).

Перетворене рівняння легко розв'язати двома послідовними інтегруваннями
(спочатку змінною ξ, а потім η). Після інтегрування на ξ отримуємо

функції \(f(x)\) у точці x у диференціальному численні.

Варіаційне числення тісно пов'язане з диференціальними рівняннями, але,
на жаль, більшість студентів його не вивчає. Ця лекція (і наступна) є
вступом до варіаційного числення та показує, як можна розв'язувати
диференціальні рівняння в похідних на основі варіаційних принципів.

Калькулюс варіацій виник одночасно з математичним аналізом у зв'язку з
розв'язанням задач максимізації та мінімізації функцій функцій (які
називаються функціоналами). Першою задачею варіаційного числення була
задача брахістохрона, сформульована Йоганном Бернуллі у 1696 році.

Бернуллі показав, що час сходження записується у формі

\begin{verbatim}
— произвольная функция от  \eta  .
\end{verbatim}

Интегрирование последнего соотношения по \textbf{η} приводит к общему
решению

і, отже, є функцією \(T\{y\}\) з функції. Оскільки багато функціоналів
впорядковані таким чином, ми зосередимося на вивченні загальної
функціональності

(44.1)

\begin{verbatim}
де  \phi(\eta)  — примітивна функція  \phi_1(\eta)  , а  \psi(\xi)  — довільна функція. Отже, загальний розв'язок рівняння
\end{verbatim}

Тепер сформульуємо головну мету наших лекцій: знайти функцію, яка дає
мінімум (або максимум) функціоналу (44.1). Стратегія пошуку буде такою
ж, як при пошуку мінімуму функції в диференціальному численні. Там ми
знайшли критичні точки функції за умовою f'(x) = 0. У варіаційному
численні все буде трохи складніше, оскільки аргумент більше не є
числовою змінною, а функцією. Однак загальний підхід залишається
незмінним: ми обчислимо так звану функціональну похідну функції \(y(x)\)
і прирівняємо її до нуля. Це нове рівняння буде аналогічним

записывается в виде

(17.3)

з диференціального числення, але тепер це буде звичайне диференціальне
рівняння, відоме як рівняння Ейлера-Лагранжа. Нам просто потрібно знайти
це рівняння і розв'язати з ним деякі проблеми.

Мінімізація функціональності

\begin{verbatim}
де  \phi(\eta)  і  \psi(\xi)  є довільними функціями своїх аргументів. Наприклад, читач може перевірити, що
\end{verbatim}

Розглянемо задачу пошуку функції \(y(x)\), яка мінімізує
функціональність

\begin{verbatim}
u(\xi, \eta) = 1/\eta + tg \xi,   u(\xi, \eta) = \eta^2 + e  являются решениями уравнения  u_{\mathrm{ }nh}=0  . Второй шаг сделан.  \text{IIIA}\Gamma  3 (Возвращение к старым координатам x и t).
\end{verbatim}

Для нахождения общего решения, т. е. всех решений уравнения
\(u_{tt} = c^2 u_{xx}\),подставим

у класі гладких функцій, які задовольняють граничні умови

, \(y(b)\) = B (див. рис. 44.2).

\begin{verbatim}
Нехай існує y = \overline{y}(x) потрібна функція, і розглянемо малу варіацію функції \overline{y} , тобто функцію \overline{y} + \varepsilon \eta(x) , де \varepsilon — мале число, а \eta(x) — гладка функція, яка задовольняє
\end{verbatim}

\begin{verbatim}
граничні умови \eta(a) = \eta(b) = 0 (див. рис. 44.2). Очевидно, що якщо обчислити інтеграл J з близької функції y + \varepsilon \eta , то функціонал буде збільшуватися, тобто
\end{verbatim}

в решение

\begin{verbatim}
за всі \epsilon . Іншими словами, графік \phi\left(\epsilon\right)=J\left[\overline{y}+\epsilon\eta\right] як функція, \epsilon буде виглядати як рис. 44.3.
\end{verbatim}

\pandocbounded{\includegraphics[keepaspectratio]{images/_page_329_Figure_5.jpeg}}

Фіг. 44.1. Завдання брахістохрани (змінна кількість завдань була
скоригована).

\pandocbounded{\includegraphics[keepaspectratio]{images/_page_329_Figure_7.jpeg}}

\begin{verbatim}
Рис. 44.2. Варіація функції: a — варіація функції \overline{y} = \overline{y}(x) + \epsilon \eta(x) ; \delta є мінімізуючою функцією \overline{y}(x) .
\end{verbatim}

З Фіг. 44.3 Очевидно, що в цьому випадку необхідно обчислити похідну від

В результате получаем

(17.4)

\begin{verbatim}
Po E, поставити \varepsilon = 0 і прирівняти отриманий вираз до нуля, тобто
\end{verbatim}

Это общее решение волнового уравнения. С физической точки эрения оно
интересно тем, что представляет сумму двух бегущих волн Р И С

\end{document}
